In this exercise, you will implement a class hierarchy that represents various types of numbers, such as rational, irrational, integer, and natural numbers. The hierarchy will be designed using inheritance, where each class extends a base class \texttt{Real}. You will also implement various mathematical operations that can be performed on these numbers.

\begin{enumerate}

    \item \textbf{Create a package named} \texttt{numbers}

    First, create a package named \texttt{numbers} where all the classes related to numbers will be stored. This helps to organize your code properly.

    \item \textbf{Implement the} \texttt{Real} \textbf{abstract class}

    The \texttt{Real} class will represent the base class for all types of numbers in the system. This class will define the common attributes and methods that all numbers share.

        \begin{itemize}

            \item \textbf{Attributes:}

                \begin{itemize}

                    \item \texttt{double value}: This attribute stores the decimal value of the real number. It will be used by subclasses to hold the value for their respective number types.

                \end{itemize}

            \item \textbf{Methods:}

                \begin{itemize}

                    \item \texttt{add(Real other)}: Adds another \texttt{Real} object to the current number and returns the result.

                    \item \texttt{subtract(Real other)}: Subtracts another \texttt{Real} object from the current number and returns the result.

                    \item \texttt{multiply(Real other)}: Multiplies the current number by another \texttt{Real} object and returns the result.

                    \item \texttt{divide(Real other)}: Divides the current number by another \texttt{Real} object and returns the result.

                \end{itemize}

        \end{itemize}

    These methods should be implemented so that they can be used by any type of number in the system. However, since \texttt{Real} is an abstract class, the subclasses are expected to implement specific behaviors where needed.

        \begin{itemize}

            \item \textbf{Constructor:}

                \begin{itemize}

                    \item \texttt{Real(double value)}: This constructor initializes the \texttt{value} attribute to a specific value passed as an argument.

                \end{itemize}

        \end{itemize}

    \textbf{Ensure that neither the constructor nor the attributes are public}. Use the appropriate visibility modifiers to restrict access but allow inheritance.

    \item \textbf{Create subclasses of} \texttt{Real}

    Now, create subclasses that represent more specific types of numbers. These classes will inherit from \texttt{Real} and implement the necessary methods for each type of number.

        \begin{itemize}

            \item \texttt{Rational} \textbf{class}: represents rational numbers, which can be expressed as a fraction.

                \begin{itemize}

                    \item \textbf{Attributes:}

                        \begin{itemize}

                            \item \texttt{int numerator}: The numerator of the rational number.

                            \item \texttt{int denominator}: The denominator of the rational number.

                        \end{itemize}

                    \item \textbf{Constructor:}

                        \begin{itemize}

                            \item \texttt{Rational(int numerator, int denominator)}: This constructor initializes both the \texttt{numerator} and \texttt{denominator} attributes. The decimal \texttt{value} is calculated internally by dividing the numerator by the denominator, and this result is passed to the parent class \texttt{Real} using \texttt{super()}.

                        \end{itemize}

                    \item \textbf{Methods:} Override the methods \texttt{add}, \texttt{subtract}, \texttt{multiply}, and \texttt{divide} to handle operations specifically using fractions if the other number is also a \texttt{Rational}, otherwise perform a standard decimal operation.

                \end{itemize}

            \item \texttt{Integer} \textbf{class}: represents integer numbers, which are a specific type of rational numbers with a denominator of 1. It should inherit directly from \texttt{Rational}.

                \begin{itemize}

                    \item \textbf{Constructor:}

                        \begin{itemize}

                            \item \texttt{Integer(int value)}: This constructor initializes the \texttt{numerator} with the given integer value and explicitly sets the \texttt{denominator} to 1 by calling the parent constructor.

                        \end{itemize}

                    \item \textbf{Methods:} Override the mathematical operation methods if necessary to ensure the results maintain integer-specific behaviors when operating with other integers.

                \end{itemize}

            \item \texttt{Natural} \textbf{class}: represents natural numbers, which are strictly positive integers used for counting. It should inherit directly from \texttt{Integer}.

                \begin{itemize}

                    \item \textbf{Constructor:}

                        \begin{itemize}

                            \item \texttt{Natural(int value)}: This constructor must ensure that the \texttt{value} is always strictly a positive integer (greater than zero). If not, it prints a warning and uses \texttt{1} instead.

                        \end{itemize}

                    \item \textbf{Methods:} Inherit mathematical operations from \texttt{Integer}, but ensure that operations resulting in negative numbers or zero are handled gracefully to maintain strict natural number rules: subtracting two naturals whose result is not positive returns an \texttt{Integer} instead of a \texttt{Natural}.

                \end{itemize}

            \item \texttt{Irrational} \textbf{class}: represents irrational numbers, such as π or the square root of 2, which cannot be expressed as exact fractions. It inherits directly from \texttt{Real}.

                \begin{itemize}

                    \item \textbf{Constructor:}

                        \begin{itemize}

                            \item \texttt{Irrational(double value)}: Initializes the \texttt{value} attribute by passing it to the parent class constructor.

                        \end{itemize}

                    \item \textbf{Methods:} The mathematical operation methods will work as expected using the inherited floating-point arithmetic, returning approximations since irrational numbers cannot be perfectly represented in a finite decimal form.

                \end{itemize}

        \end{itemize}

    \item \textbf{Main program}

    Create a \texttt{Main} class that will test all the implemented classes and their functionality.

        \begin{itemize}

            \item Instantiate objects of \texttt{Integer}, \texttt{Natural}, \texttt{Rational}, and \texttt{Irrational}.

            \item Perform arithmetic operations like addition, subtraction, multiplication, and division on these objects and print the results to demonstrate polymorphism in action.

        \end{itemize}

\end{enumerate}
